%% notes-tex/database-expansion-bacteria/sections/rule.tex
%% [[experiments.database.expansion-bacteria]]
%% https://github.com/Mjvolk3/torchcell/tree/main/notes-tex/database-expansion-bacteria/sections/rule.tex
%% SOURCE: hand-authored prose; every number quoted here comes from
%% experiments/database/scripts/build_bacteria_candidate_datasets_table.py

\section{The selection rule\secstatus{todo}}\label{sec:rule}

\subsection{Rows are ordered by measurements, which is scale times density\secstatus{todo}}

The order is one quantity, applied to every row alike: \term{measurements}, meaning
instances times phenotype dimensionality, descending. \term{Instances} is the number of
genotype $\times$ environment $\times$ replicate records a row can actually yield, and
\term{dimensionality} is how many values each of those records carries. A scalar fitness
screen of $n$ strains contributes $n$ numbers; a proteome panel of $n$ samples at a
depth of 2,000 proteins contributes $2\times10^{3}n$.

Ranking on instances alone gets the order visibly wrong, and it is worth naming the
failure rather than asserting the rule. A quantitative proteome across a couple of dozen
growth conditions holds fewer instances than a small deletion screen, which would rank it
below one, and yet each of its instances is a complete proteome. The product is what
places it correctly, and it is the same normalization the supported-datasets table
already applies through its compressed-signal column and the same one the yeast candidate
list ranks inside its bands with.

Two tie-breaks make the rule total. Instances break a measurement tie, so between two
rows holding the same number of values the one spanning more distinct strain and
condition combinations wins; a wide, shallow row is worth more to a
genotype-to-phenotype model than a narrow, deep one of equal volume. Tier breaks that.

This differs from the yeast list on purpose. There, rows are banded first by what they
supply to a planned single-cell campaign, and scale orders only inside a band. Here scale
and density are the requested rule directly, so no band is applied and no row is promoted
for what it pairs with. What a row combines with is still recorded, in
Table~\ref{tab:banalogs} and in the join column, but it does not move the row.

\subsection{Sequence reconstruction is the gate, not a preference\secstatus{todo}}

A phenotype is only trainable against a genotype that can be written down, so every row
names a \term{sequence basis}: the route by which the total genomic content of one strain
would be reconstructed. The bases are named per host, because the reference differs and
because a b-number and a PP\_ locus tag are different namespaces that must not be
silently merged (Sec.~\ref{sec:schema}).

\begin{itemize}[leftmargin=1.2em,itemsep=1pt,topsep=2pt]
  \item \term{MG1655-KO} and \term{KT2440-KO} -- the reference genome minus one
    cataloged gene. The cheapest basis, and the one most of the built yeast set sits on
    in its own reference.
  \item \term{MG1655-KO x KO} -- a constructed double mutant, both loci cataloged, which
    is the basis a genetic-interaction row needs.
  \item \term{MG1655+transposon} and \term{KT2440+transposon} -- a mapped insertion site
    carrying a barcode. The insertion position is sequenced, so the genotype is realized
    rather than designed, which is what separates a barcoded transposon library from a
    guide library.
  \item \term{MG1655+guide}, \term{KT2440+guide} -- the genome is unedited and the
    perturbation is a designed cassette plus a guide target. Regulatory, not
    sequence-level, and it reaches the essential genes a deletion collection cannot hold.
  \item \term{KT2440+promoter} and \term{engineered-chassis+promoter} or
    \term{+RBS} -- a designed regulatory part over a native or heterologous gene.
    Designed rather than verified per strain.
  \item \term{engineered-chassis} -- a named production strain plus heterologous
    cassettes, which is what a titer is measured on.
  \item \term{evolved-WGS} -- a resequenced evolved clone.
  \item \term{reference-only} -- wild type, where the environment carries the
    perturbation.
\end{itemize}

What fails the gate is a strain whose genome is unknown even in principle. An evolved
population that was never resequenced lands there, and so does a strain described only by
a phenotype in a review table. Those are dropped in Table~\ref{tab:bexcluded} rather than
ranked low, because no amount of phenotypic value substitutes for a missing genotype.

\subsection{Tier, and what it is for here\secstatus{todo}}

Tier no longer orders the list, so it does one job: it records how much of a dataset a
loader would actually get, which is what makes a high measurement count trustworthy or
not.

\begin{itemize}[leftmargin=1.2em,itemsep=1pt,topsep=2pt]
  \item \textbf{Tier 1} clears all three of $\geq$$10^3$ genotypes, $\geq$$10^4$
    instances, and a clean sequence basis.
  \item \textbf{Tier 2} clears two of the three, or is the only dataset covering a
    phenotype the substrate needs for these hosts.
  \item \textbf{Tier 3} is a modality or reference backbone the substrate needs without
    itself carrying a large genotype axis, such as a library definition or a per-gene
    coefficient that applies to every strain alike.
  \item \textbf{Tier 4} is real but small, coarsely labeled, or expensive to extract.
\end{itemize}

The bars apply to what a row contains rather than to its reputation. A curated corpus
with thousands of literature records fails tier 1 whatever its size, because its records
re-serve primary papers and the net-new fraction is unknown until it is split by source.

\subsection{Staging, and why the first tranche carries a quota\secstatus{todo}}\label{sec:stage}

The fifty are built in two tranches, twenty then thirty. The second tranche is the next
thirty by the ordering rule and needs no further explanation. The first twenty carry one
correction to the rule, and it is applied after ranking and reported rather than folded
into a score.

Measurements alone would let one host fill the first tranche, because the two literatures
are not the same size and the largest screens are not evenly split between them. Starting
both hosts together is worth more than starting the twenty largest rows: the two schema
changes in Sec.~\ref{sec:schema} are per host, the provisioning of a reference release is
per host, and neither cost is paid until a row for that host is actually built. So the
first tranche holds ten rows per host, taken in rank order within each, and every row the
quota promoted is named with its unaided rank in Table~\ref{tab:bfinal} and in the
generating script's output.

The consequence is visible and intended. A row promoted by the quota sits above rows with
more measurements than it has, and the displaced rows are the weakest of the
over-represented host rather than the weakest overall, so a promotion never costs a row
the quota is protecting.
